\documentclass[11pt,a4paper]{article}
\usepackage[T1]{fontenc}
\usepackage[utf8]{inputenc}
\usepackage{lmodern}
\usepackage[margin=2.5cm]{geometry}
\usepackage{booktabs}
\usepackage{tabularx}
\usepackage{longtable}
\usepackage{enumitem}
\usepackage{listings}
\usepackage{xcolor}
\usepackage[hidelinks]{hyperref}

\lstset{
  basicstyle=\ttfamily\small,
  columns=fullflexible,
  keepspaces=true,
  frame=single,
  rulecolor=\color{gray!50},
  backgroundcolor=\color{gray!8},
  breaklines=true,
  upquote=true,
  xleftmargin=0.5em,
  xrightmargin=0.5em
}
\setlength{\emergencystretch}{3em}
\newcommand{\kbd}[1]{\texttt{#1}}
\newcolumntype{L}{>{\raggedright\arraybackslash}X}

\title{LeRobot User Manual}
\author{}
\date{9 November 2026}

\begin{document}
\maketitle

LeRobot is a drawing robot you steer on screen: either live, with the mouse
and keyboard (Dessiner), or by writing a short OCaml program of robot
commands (such as \kbd{mon\_dessin.ml}). It is the OCaml translation of the
1998 Caml Light version, in the folder \kbd{LeRobot\_OCaml\_FR} of the
\kbd{LeRobot} folder, which can be anywhere (commands below are typed in a
Terminal window opened in the \kbd{LeRobot} folder).

\tableofcontents

%======================================================================
\section{Getting started}

You need OCaml with dune and the \kbd{graphics} library
(\kbd{opam install graphics}), and XQuartz, the X11 display server for macOS.

\begin{enumerate}
\item Build once, and again after any change to the source files:
\begin{lstlisting}[language=bash]
cd LeRobot_OCaml_FR && eval $(opam env) && dune build
\end{lstlisting}
\item Start XQuartz: \kbd{open -a XQuartz}. Opening the robot window starts
  it too if it is set up correctly.
\item Always launch with \kbd{DISPLAY=:0} in front of the command, as shown in
  Parts~1 and~3 (\kbd{robot.sh} and \kbd{LeRobotOCaml.app} do it for you). The default display setting can point to a stale XQuartz
  connection and fail with \emph{Cannot open display}.
\end{enumerate}

The application \kbd{LeRobotOCaml.app}, in \kbd{LeRobot\_OCaml\_FR}, does all
of this for you: double-click it to run Dessiner, or drop a \kbd{.ml}
program on its icon to compile and run it (see Part~2).

After a restart of the Mac, XQuartz may refuse to start because the folder
\kbd{/tmp/.X11-unix} is missing or not owned by root. The startup task
described in the installation manual recreates it at every restart; without
it, recreate it once (asks for your password):
\begin{lstlisting}[language=bash]
sudo mkdir -p /tmp/.X11-unix && sudo chown root:wheel /tmp/.X11-unix && sudo chmod 1777 /tmp/.X11-unix
\end{lstlisting}

The window is $800 \times 700$ pixels. The large white area is where the
robot draws; the strip along the bottom is the message bar.

%======================================================================
\section{Part 1 --- Interactive drawing (Dessiner)}

In Dessiner you give the robot commands with single keys and move it with
the mouse or the arrow keys; press Enter when you are done, then \kbd{y} to
save the drawing as a program.

\subsection{Launch}
Double-click \kbd{LeRobotOCaml.app}, or type \kbd{make Dessiner}, or:
\begin{lstlisting}[language=bash]
cd LeRobot_OCaml_FR && DISPLAY=:0 dune exec ./Dessiner.exe
\end{lstlisting}
The robot appears in the centre, pointing north, as a small arrow-shaped
icon. To give the window keyboard focus, click its title bar.

\subsection{The window}
\begin{itemize}
\item \textbf{Drawing area} --- the large white area.
\item \textbf{Message bar} (bottom) --- a reminder of the keys, ending with
  \kbd{* :cc w :peindre\ \ Entree :fin}.
\item \textbf{Colour box} (small box near the right of the bar) --- shows
  the current pen colour.
\item \textbf{Enter icon} (far right) --- the Enter symbol; it blinks when
  the robot waits for a key.
\end{itemize}

\subsection{Moving the robot with the mouse or the arrow keys}
Click on a grid point (or move the pointer with the button held down) and
the robot jumps to the nearest grid point, without drawing. The arrow keys
move it one grid step left, right, up or down, also without drawing. When
the pointer moves with the button released, the robot does not move, and a
click or trackpad tap does not end the drawing.

If the robot would go past the edge of the drawing area, it stays where it
is and the terminal prints
\kbd{** SORTIE INTERDITE ** translater} (``forbidden exit''), with a beep-beep.
This is only a warning.

\subsection{Keyboard commands}
The robot draws only when its pen is down (\kbd{b}). One step is one grid
square.

\begin{center}
\begin{tabularx}{\linewidth}{@{}l l L@{}}
\toprule
\textbf{Key} & \textbf{Command} & \textbf{What it does} \\
\midrule
\kbd{b} & bc & Pen down: the robot draws as it moves \\
\kbd{l} & lc & Pen up: the robot moves without drawing \\
\kbd{a} & av & Move forward one step \\
\kbd{)} & vd & Curved turn to the right (draws an arc) \\
\kbd{(} & vg & Curved turn to the left (draws an arc) \\
\kbd{+} & pqtd & Quarter turn right ($90^\circ$) \\
\kbd{-} & pqtg & Quarter turn left ($90^\circ$) \\
\kbd{d} & phtd & Eighth turn right ($45^\circ$) \\
\kbd{g} & phtg & Eighth turn left ($45^\circ$) \\
\kbd{c} & dc & Draw a small cross (+) \\
\kbd{p} & dp & Draw a point \\
\kbd{.} & dg & Draw the grid \\
\kbd{*} & cc & Next pen colour (shown in the colour box) \\
\kbd{w} & peindre & Paint the closed area behind the robot with the pen colour \\
Arrow keys & translater & Move one step left, right, up or down, without drawing \\
Backspace & --- & Undo the last command \\
Enter & --- & Finish the drawing \\
\bottomrule
\end{tabularx}
\end{center}

Colours cycle in this order: white, yellow, blue, violet, green, red,
indigo, black. To paint, put the robot so that its back is inside a closed
area, then press \kbd{w}.

\subsection{Finishing and saving}
\begin{enumerate}
\item Press \textbf{Enter}. The bar asks:
  \emph{Voulez-vous sauver dans mon\_dessin.ml ? (y/n)}
\item Press \kbd{y} to save. Every command you used, including mouse and
  arrow moves (written as \kbd{translater (dx,dy)}) and colour changes
  (written as \kbd{cc n} with the colour name in a comment), is written to
  \kbd{mon\_dessin.ml} as a robot program (see Part~3), and the window
  closes.
\item Press any other key to close without saving.
\end{enumerate}
Saving always overwrites \kbd{mon\_dessin.ml}. To keep an earlier drawing,
copy that file under another name first.

%======================================================================
\section{Part 2 --- Running a robot program}
A robot program is an OCaml file that lists robot commands in order. It can
be in any folder. Drop it on \kbd{LeRobotOCaml.app}, or type, in a Terminal
window opened in \kbd{LeRobot\_OCaml\_FR}:
\begin{lstlisting}[language=bash]
./robot.sh path/to/Program.ml
\end{lstlisting}
\kbd{robot.sh} compiles the program (the library must have been built once
with \kbd{make}), starts XQuartz if needed and runs it. The program draws the
picture by itself, then waits so you can replay it at another speed (see
Section~\ref{sec:replay}) or close it with Esc.

The programs of the book \emph{Premi\`eres le\c{c}ons de programmation en
OCaml} start with \kbd{open LibRobot} and are in the folder \kbd{livre},
one subfolder per chapter:
\begin{lstlisting}[language=bash]
./robot.sh livre/03_programmation/Carre.ml
make run PROG=livre/20_grille_polaire/LimaconDePascal
make verifier                      # all programs, one after the other
make verifier DIR=02_interaction   # one chapter
make verifier DEPUIS=FriseDeTriangles   # resume at a program
\end{lstlisting}
With \kbd{make verifier}, answer in the robot window when a drawing is
finished: Enter = correct (next program), \kbd{r} = run it again, \kbd{q}
or Esc = stop (the command to resume is printed), any other key starts a
remark, ended by Enter, noted in \kbd{verification.txt}. A program that
never ends by itself (\kbd{Boucle1}, \kbd{Boucle2}) is stopped with Esc.

%======================================================================
\section{Part 3 --- Running a robot program with dune}

A robot program is an OCaml file that lists robot commands in order; you run
it with \kbd{dune exec}, and it draws the picture by itself, then waits so
you can replay it at another speed.

\subsection{Run mon\_dessin.ml}
\begin{lstlisting}[language=bash]
cd LeRobot_OCaml_FR && DISPLAY=:0 dune exec ./mon_dessin.exe
\end{lstlisting}
\kbd{dune exec} rebuilds the program first if the file has changed.

\subsection{Write your own program}
\begin{enumerate}
\item Create a file in \kbd{LeRobot\_OCaml\_FR}, for example \kbd{maison.ml}:
\begin{lstlisting}
open LeRobot
let () =
begin
  init_pict (); ce();
  bc(); av(); av(); pqtd(); av(); av(); pqtd(); av(); av(); pqtd(); av(); av();
  st()
end
\end{lstlisting}
\item Add its name, without \kbd{.ml}, to both lists of the
  \kbd{executables} section in the file \kbd{dune}:
\begin{lstlisting}
(executables
 (names Dessiner mon_dessin maison)
 (modules Dessiner mon_dessin maison)
\end{lstlisting}
\item Run it: \kbd{DISPLAY=:0 dune exec ./maison.exe} from the
  \kbd{LeRobot\_OCaml\_FR} folder.
\end{enumerate}

Every program follows the same pattern:
\begin{itemize}
\item \kbd{init\_pict (); ce();} first: clears the window and puts the robot
  in the centre.
\item The robot starts with its pen down, so every move draws: call
  \kbd{lc()} to move without drawing, and \kbd{bc()} to draw again.
\item \kbd{st()} last: shows the speed screen and keeps the window open.
  Without it, the window closes as soon as the drawing is done.
\item Each command is followed by \kbd{;}, and commands without arguments
  take \kbd{()}.
\end{itemize}

\subsection{Replaying at another speed}\label{sec:replay}
When the program reaches \kbd{st()}, the bar shows:
\begin{lstlisting}
+ : plus vite   - : moins vite   vitesse : N     autre touche : redessiner   Echap : quitter
\end{lstlisting}
\begin{itemize}
\item \kbd{+} / \kbd{-} --- change the speed from 0 to 10 (default 10, the
  fastest). Only the number changes.
\item Any other key, for example the space bar --- clear the window and
  redraw at the chosen speed; the speed screen comes back when it is done.
\item Esc --- close the window.
\end{itemize}
At speed 1 each command takes about 0.1\,s. Speed 0 is step-by-step mode:
the robot does one command each time you press a key, and Esc stops.

\subsection{Command reference}
One step is one grid square. Commands are in French; the table gives the
English meaning.

\begin{longtable}{@{}p{0.42\linewidth} p{0.53\linewidth}@{}}
\toprule
\textbf{Command} & \textbf{Meaning} \\
\midrule
\endhead
\bottomrule
\endfoot
\kbd{av ()} & Move forward one step \\
\kbd{avf n} & Move forward $n$ steps \\
\kbd{vd ()} / \kbd{vg ()} & Curved turn right / left (draws an arc) \\
\kbd{phtd ()} / \kbd{phtg ()} & Turn $45^\circ$ right / left \\
\kbd{pqtd ()} / \kbd{pqtg ()} & Turn $90^\circ$ right / left \\
\kbd{p3htd ()} / \kbd{p3htg ()} & Turn $135^\circ$ right / left \\
\kbd{pdt ()} & Turn around ($180^\circ$) \\
\kbd{pn ()} \kbd{pne ()} \kbd{pe ()} \kbd{pse ()} \kbd{ps ()} \kbd{psw ()} \kbd{pw ()} \kbd{pnw ()} & Face N, NE, E, SE, S, SW, W, NW \\
\kbd{pivoter n} & Face direction $n$ (0 = N, 1 = NE, 2 = E \dots\ 7 = NW) \\
\kbd{rt a} & Turn by $a$ degrees (anticlockwise) \\
\kbd{bc ()} / \kbd{lc ()} & Pen down / pen up \\
\kbd{dp ()} / \kbd{dc ()} & Draw a point / a small cross \\
\kbd{cc n} & Pen colour $n$ (0 white, 1 yellow, 2 blue, 3 violet, 4 green, 5 red, 6 indigo, 7 black) \\
\kbd{ec n} & Pen thickness $n$ \\
\kbd{cf n} & Clear the window and fill it with background colour $n$ \\
\kbd{translater (dx,dy)} & Move $dx$ steps across and $dy$ steps up, without drawing \\
\kbd{tracer (x,y)} & Go to window pixel $(x,y)$, drawing if the pen is down \\
\kbd{deplacer (x,y)} & Go to grid point $(x,y)$, counted from the centre \\
\kbd{ce ()} & Go to the centre \\
\kbd{eh ()} \kbd{eb ()} \kbd{ad ()} \kbd{ag ()} & Go to the top, bottom, right or left edge \\
\kbd{dg ()} & Draw the grid \\
\kbd{lg x} / \kbd{lgX x} / \kbd{lgY x} & Set the step size (both directions / across / up) \\
\kbd{lgRT (r,t)} & Step of length $r$ at angle $t$ (polar grid) \\
\kbd{tl l} & Move forward $l$ points, the grid following the robot \\
\kbd{peindre ()} & Fill the closed area behind the robot with the pen colour \\
\kbd{montrerRobot ()} / \kbd{cacherRobot ()} & Show / hide the robot \\
\kbd{autoriserSorties ()} / \kbd{interdireSorties ()} & Allow / forbid moving past the window edge \\
\kbd{vt n} & Set the speed (0 to 10) \\
\kbd{pp ()} / \kbd{ex ()} & Step-by-step mode / full speed \\
\kbd{message "text"} & Write a message in the bottom bar \\
\kbd{st ()} & End: show the speed screen \\
\end{longtable}

The full list, with test functions such as \kbd{tb ()} (would the next step
leave the window?), is in \kbd{LeRobot.mli}.

%======================================================================
\section{Troubleshooting}

\begin{longtable}{@{}p{0.27\linewidth} p{0.33\linewidth} p{0.34\linewidth}@{}}
\toprule
\textbf{Symptom} & \textbf{Cause} & \textbf{Fix} \\
\midrule
\endhead
\bottomrule
\endfoot
\emph{Cannot open display} & DISPLAY points to a stale XQuartz connection, or XQuartz is not running & Put \kbd{DISPLAY=:0} in front of the command; start XQuartz with \kbd{open -a XQuartz} \\
\addlinespace
XQuartz opens and closes at once & \kbd{/tmp/.X11-unix} missing or not owned by root (common after a restart) & Run the \kbd{sudo mkdir \dots} command from Section~1 \\
\addlinespace
\kbd{** SORTIE INTERDITE ** translater} in the terminal & The robot was asked to move past the window edge & Only a warning; keep the pointer or the robot inside the window, or call \kbd{autoriserSorties ()} \\
\addlinespace
The window closes right after the drawing & The program has no \kbd{st()} at the end & Add \kbd{st()} as the last command \\
\addlinespace
Keys do nothing in the window & The window does not have keyboard focus & Click the window's title bar \\
\addlinespace
A saved drawing was lost & Saving from Dessiner always overwrites \kbd{mon\_dessin.ml} & Copy \kbd{mon\_dessin.ml} under a new name before saving again \\
\addlinespace
Text font cannot be changed & \kbd{policeTexte} has no effect with this XQuartz & None; the bar layout is sized for the default font \\
\end{longtable}

\end{document}
